%%=============================================================================
%% Samenvatting
%%=============================================================================

% TODO: De "abstract" of samenvatting is een kernachtige (~ 1 blz. voor een
% thesis) synthese van het document.
%
% Een goede abstract biedt een kernachtig antwoord op volgende vragen:
%
% 1. Waarover gaat de bachelorproef?
% 2. Waarom heb je er over geschreven?
% 3. Hoe heb je het onderzoek uitgevoerd?
% 4. Wat waren de resultaten? Wat blijkt uit je onderzoek?
% 5. Wat betekenen je resultaten? Wat is de relevantie voor het werkveld?
%
% Daarom bestaat een abstract uit volgende componenten:
%
% - inleiding + kaderen thema
% - probleemstelling
% - (centrale) onderzoeksvraag
% - onderzoeksdoelstelling
% - methodologie
% - resultaten (beperk tot de belangrijkste, relevant voor de onderzoeksvraag)
% - conclusies, aanbevelingen, beperkingen
%
% LET OP! Een samenvatting is GEEN voorwoord!

%%---------- Nederlandse samenvatting -----------------------------------------
%
% TODO: Als je je bachelorproef in het Engels schrijft, moet je eerst een
% Nederlandse samenvatting invoegen. Haal daarvoor onderstaande code uit
% commentaar.
% Wie zijn bachelorproef in het Nederlands schrijft, kan dit negeren, de inhoud
% wordt niet in het document ingevoegd.

\IfLanguageName{english}{%
\selectlanguage{dutch}
\chapter*{Samenvatting}
\lipsum[1-4]
\selectlanguage{english}
}{}

%%---------- Samenvatting -----------------------------------------------------
% De samenvatting in de hoofdtaal van het document

\chapter*{\IfLanguageName{dutch}{Samenvatting}{Abstract}}

De toegankelijkheid van huisartsenpraktijken staat in Vlaanderen onder druk. 
Een vergrijzende artsenpopulatie, een dalende instroom in de huisartsopleiding en een stijgende zorgvraag zorgen voor een structurele belasting van de eerstelijnszorg.
Om deze problematiek te monitoren, ontwikkelden organisaties zoals Domus Medica en UGent gevalideerde indicatorensets en vragenlijsten. 
Uit pilootstudies blijkt echter dat huisartsen en praktijkmedewerkers deze evaluatie-instrumenten als te tijdrovend, complex en inconsistent ervaren, wat de kwaliteit en frequentie van de dataverzameling ernstig beperkt.

Deze bachelorproef onderzoekt in welke mate een digitale proof-of-concept interface de toegankelijkheidsevaluatie voor huisartsen en praktijkmedewerkers sneller, consistenter en gebruiksvriendelijker kan maken dan de huidige aanpak. 
Het onderzoek werd uitgevoerd binnen de context van Eerstelijnszone Aalst en vertrekt van de gevalideerde indicatorenset uit het rapport \textit{Toegankelijkheid Huisartsgeneeskunde} \autocite{Merckx2023}.

Het onderzoek verliep in vier fasen. 
Eerst werd via een werkgroep en literatuurstudie het ontwerp van de interface bepaald en gevalideerd. 
Vervolgens werd de interface ontwikkeld met Angular als frontend-framework, Express.js als backend en PostgreSQL als databank. 
SurveyJS werd ingezet voor de vragenlijstcomponent, Auth0 voor authenticatie en Microsoft Clarity voor gedragsanalyse. 
In een derde fase werd de interface getest door twaalf huisartsen en praktijkmedewerkers. 
Tot slot werden de resultaten kwantitatief en kwalitatief geanalyseerd.

De gebruikerstesten tonen aan dat de interface als goed bruikbaar wordt ervaren, met een gemiddelde SUS-score van 73,96. 
De lage scores op de negatieve stellingen bevestigen dat de applicatie een lage instapdrempel heeft en zonder technische kennis zelfstandig kan worden ingevuld. 
De tijdsmetingen via Microsoft Clarity tonen dat de eigenlijke vragenlijst een mediane invultijd van 2 minuten en 28 seconden kent. 
De kwalitatieve feedback bevestigt de inhoudelijke meerwaarde van de interface, maar wijst ook op verbeterpunten zoals het ontbreken van automatisch opslaan en inconsistenties in de interface-elementen.

De bachelorproef toont aan dat een gerichte digitale interface het evaluatieproces voor huisartsen concreet kan vereenvoudigen. 
De resultaten bieden een basis voor verdere ontwikkeling, onder meer via een koppeling met bestaande datasystemen zoals RIZIV-registers en de integratie van gepersonaliseerde rapportgeneratie via een taalmodel.
